\chapter{Implementation of the Distributed Lakehouse Platform}

\section{Introduction}
This chapter presents the concrete implementation of \projectname{} as an operational distributed lakehouse platform. While Chapter 4 defined the design baseline, the present chapter explains how that design was translated into working services, executable pipelines, governed data layers, and secured access paths. The objective is not to describe isolated tools separately, but to show how ingestion, storage, processing, metadata, querying, transformation, quality validation, and orchestration were assembled into one coherent engineering system.

\section{Data Sources and Dataset Modeling}
The implementation is based on a compact but representative dataset set selected to support the Customer 360 orientation of the platform while remaining compatible with the available academic infrastructure. Rather than using one oversized monolithic source, the project retains several complementary datasets, each contributing a different view of customer behavior.

\begin{table}[!htbp]
    \centering
    \footnotesize
    \renewcommand{\arraystretch}{1.16}
    \arrayrulecolor{tablelineblue}
    \rowcolors{2}{tablebodyblue}{white}
    \begin{tabularx}{\textwidth}{|Y|Y|Y|Y|}
        \hline
        \rowcolor{tableheaderblue}
        \textcolor{white}{\textbf{Dataset}} & \textcolor{white}{\textbf{Analytical role}} & \textcolor{white}{\textbf{Customer signal covered}} & \textcolor{white}{\textbf{Implementation contribution}} \\
        \hline
        Bank Marketing & campaign and response analysis & contact history, subscription outcome, profile attributes & provides customer profile and campaign-oriented raw inputs \\
        \hline
        Online Retail II & transactional and commercial behavior analysis & invoices, products, quantities, revenue, countries & provides purchase history and product interaction facts \\
        \hline
        Online Shoppers Purchasing Intention & digital behavioral analysis & session behavior, navigation, purchase intention & provides pre-purchase intent and online behavior indicators \\
        \hline
    \end{tabularx}
    \caption{Active datasets retained in the implemented platform}
    \label{tab:chapter5-datasets}
\end{table}

Each dataset keeps its source identity during ingestion and bronze persistence. This choice is important because the platform is intended to preserve replayability and traceability before any normalization or business transformation is applied. In practice, the implementation therefore distinguishes between source-specific raw inputs and later converged analytical models. This allows the platform to defend both schema fidelity at the entry point and business-oriented integration in upper layers.

\textit{Diagram placeholder: active dataset families and how each one contributes to the customer-oriented analytical scope of the platform.}
\FloatBarrier

\section{Kafka-Based Ingestion Layer}
Kafka is used as the transport backbone of the ingestion layer. The implementation does not load datasets directly into storage or processing services; instead, each source is first published through a topic-oriented flow that decouples dataset reading from downstream persistence. This approach improves control, replayability, and operational clarity.

The Kafka cluster is distributed across the three virtual machines:
\begin{itemize}
    \item broker 1 on VM1;
    \item broker 2 on VM2;
    \item broker 3 on VM3.
\end{itemize}

At the implementation level, each source dataset is associated with its own topic family. This prevents unrelated sources from being merged too early and makes reset and replay operations more explicit. The control plane triggers the ingestion logic, producers publish the records, and downstream consumers persist the bronze output in HDFS. This means Kafka acts as the controlled transport boundary between local source preparation and distributed storage.

The platform also implements topic reset logic within the orchestration flow. This is important for demonstrations and reproducible reruns, because it allows the environment reset pipeline to clean old ingestion traces before a fresh execution begins.

\textit{Diagram placeholder: dataset publication flow from the control plane to the three-broker Kafka cluster and toward bronze persistence.}

\section{HDFS Bronze Storage Layer}
The bronze layer is implemented on HDFS as the replayable raw persistence zone of the platform. After Kafka-based transport, source records are landed into HDFS in batch-oriented bronze files before any lakehouse structuring step is executed. This choice gives the project an explicit raw persistence boundary, which is essential for reproducibility, debugging, and future reprocessing.

The HDFS cluster is implemented with:
\begin{itemize}
    \item the NameNode on VM2;
    \item one DataNode on each VM: VM1, VM2, and VM3.
\end{itemize}

In operational terms, this means the namespace authority remains centralized on VM2, while block storage is distributed across the three nodes. The implemented bronze organization is client-oriented and source-aware: each dataset is persisted under a controlled bronze structure rather than being mixed into one generic folder. This reinforces traceability and keeps the raw layer easy to inspect.

The bronze layer also plays a strategic role in the whole project: it is the stable checkpoint between ingestion and compute. If a downstream Spark stage fails, the platform does not need to reread the original local files or republish the full datasets manually. It can simply restart from the bronze persistence already distributed in HDFS.

\textit{Diagram placeholder: bronze storage organization in HDFS, showing NameNode coordination and DataNode distribution across VM1, VM2, and VM3.}

\section{Spark Processing and Iceberg Loading}
Spark is the main distributed processing engine of the platform. Its implementation is aligned with the physical cluster design:
\begin{itemize}
    \item the Spark master runs on VM2;
    \item workers are distributed across VM1, VM2, and VM3;
    \item the control plane submits jobs remotely to the Spark environment on VM2.
\end{itemize}

This execution path is important. The final deployment does not rely on a heavy local Spark driver on the workstation. Instead, the control plane orchestrates remote execution through the distributed cluster, which is closer to a production-style separation between orchestration and compute.

From a data perspective, Spark reads the bronze files from HDFS, interprets their schemas, normalizes the raw structures, and writes the outputs into Iceberg-backed raw lakehouse tables. The implementation therefore transforms replayable raw artifacts into governed table objects while preserving distributed execution. The raw lakehouse layer is not yet the final curated business layer; it is the structured foundation on top of which transformations and validations are later applied.

\textit{Diagram placeholder: Spark raw-load path from HDFS bronze files to Iceberg raw tables through the distributed Spark cluster.}

\section{Hive Metastore and Catalog Governance}
Hive Metastore is implemented as the shared metadata authority of the lakehouse platform. It is deployed on VM2 together with its PostgreSQL backing store and acts as the central catalog reference for upper services such as Spark and Trino.

Its role in the implementation is essential: once Spark writes the raw Iceberg tables, their definitions must be discoverable through a common metadata layer. Without this shared governance service, each engine would have to manage table visibility independently, which would weaken consistency and make the architecture harder to defend.

In practical terms, Hive Metastore allows the platform to:
\begin{itemize}
    \item register the lakehouse tables created after Spark loading;
    \item expose the same logical objects to different access engines;
    \item preserve a governed catalog boundary between raw persistence and analytical consumption.
\end{itemize}

This makes the metadata layer a real implementation component, not a theoretical addition. It is the reason the lakehouse can behave like a coordinated data platform rather than a simple collection of files.

\textit{Diagram placeholder: central metadata governance through Hive Metastore between Spark writers and Trino readers.}

\section{Trino Query Layer}
Trino provides the analytical SQL access layer of the platform. Its implementation follows a coordinator/worker separation:
\begin{itemize}
    \item the coordinator is hosted on VM2;
    \item worker 1 is hosted on VM1;
    \item worker 2 is hosted on VM3.
\end{itemize}

This arrangement is important because it demonstrates that the query layer is also distributed, not only the storage or compute layer. Trino consumes the shared catalog definitions and exposes interactive SQL access over the Iceberg tables. In the project, this layer is used for readiness checks, lakehouse verification, dataset inspection, and analytical exploration.

The query access path is also protected by an authenticated gateway. As a result, the Trino browser interface is not exposed as an uncontrolled plain endpoint. This reinforces the idea that access services in the final platform must remain demonstrable but still aligned with the security posture of the deployment.

\textit{Diagram placeholder: Trino coordinator/worker query path and the secure gateway used for analytical access.}

\section{dbt-spark Transformation Layer}
dbt-spark is used to organize transformations as declarative models rather than as ad hoc processing scripts. This implementation choice is important because it introduces structure, dependency control, and documentation logic into the transformation stage.

The implemented transformation flow follows a layered modeling philosophy:
\begin{itemize}
    \item raw lakehouse tables remain the direct output of the Spark loading stage;
    \item staging models standardize and clean the raw structures;
    \item intermediate models prepare reusable analytical logic;
    \item analytical models expose business-oriented outputs ready for exploration.
\end{itemize}

In this design, dbt does not replace Spark. Instead, it builds on top of Spark Thrift access and uses the distributed lakehouse environment as its execution target. This distinction matters in the report because it shows that transformation governance is separated from raw compute mechanics: Spark provides the engine, while dbt provides the modeling discipline.

\textit{Diagram placeholder: layered dbt transformation flow from raw tables to staging, intermediate, and analytical outputs.}

\section{Great Expectations Validation Layer}
Great Expectations is implemented as the lightweight validation layer of the platform. Its role is not to perform generic visual checking, but to formalize trust conditions over the produced data. This is especially important in an academic lakehouse platform, where the project must prove that outputs are not only generated, but also controlled.

The implemented validation logic targets the structural and semantic integrity of selected lakehouse layers. In practice, the platform uses Great Expectations to check that the resulting data products remain consistent with expected shapes and business-level constraints. Validation artifacts are then persisted so that pipeline runs can be defended with observable evidence instead of informal claims.

This layer is therefore central to the governance dimension of the project. It strengthens the transition from “data processed successfully” to “data can be considered trustworthy enough for analytical use.”

\textit{Diagram placeholder: validation checkpoints over the lakehouse layers and persistence of validation evidence.}

\section{Airflow Orchestration Layer}
Airflow orchestrates the platform from the local control plane. In the final implementation, orchestration is not reduced to one oversized DAG. Instead, the project is structured around a small set of operational DAGs, each with a clear responsibility in the lifecycle of the environment.

\begin{table}[!htbp]
    \centering
    \footnotesize
    \renewcommand{\arraystretch}{1.16}
    \arrayrulecolor{tablelineblue}
    \rowcolors{2}{tablebodyblue}{white}
    \begin{tabularx}{\textwidth}{|Y|Y|Y|}
        \hline
        \rowcolor{tableheaderblue}
        \textcolor{white}{\textbf{DAG}} & \textcolor{white}{\textbf{Operational role}} & \textcolor{white}{\textbf{Implementation value}} \\
        \hline
        Environment reset pipeline & cleans Kafka, HDFS bronze traces, and lakehouse state & enables reproducible reruns from a controlled baseline \\
        \hline
        Lakehouse setup pipeline & prepares technical structures and service prerequisites & ensures the catalog and storage environment are initialized consistently \\
        \hline
        Lakehouse readiness pipeline & validates that the main services are reachable and usable & prevents execution from starting on an unstable cluster \\
        \hline
        Kafka--HDFS--Spark lakehouse pipeline & executes the raw ingestion and loading flow & materializes the raw distributed lakehouse foundation \\
        \hline
        dbt-spark lakehouse pipeline & runs modeled transformations and validation-oriented upper steps & produces curated analytical outputs on top of the raw layer \\
        \hline
    \end{tabularx}
    \caption{Implemented orchestration flow of the platform}
    \label{tab:chapter5-dags}
\end{table}

This orchestration strategy is important because it matches the real operational logic of the project. The platform can be reset, checked, loaded, transformed, and validated in separate controlled stages. It also makes failures easier to diagnose, because a problem can be linked to a specific orchestration boundary instead of being hidden inside one monolithic workflow.

\textit{Diagram placeholder: five-DAG orchestration sequence from environment preparation to raw loading and curated transformation.}
\FloatBarrier

\section{Conclusion}
This chapter showed how the design baseline of the previous chapter was turned into a concrete distributed implementation. The platform now appears as a full engineering chain: heterogeneous customer datasets are published through Kafka, persisted in HDFS bronze storage, processed through Spark, structured into Iceberg tables, governed through Hive Metastore, queried through Trino, transformed through dbt-spark, validated through Great Expectations, and orchestrated through Airflow.

The implementation therefore confirms that the project is not a simple tool collection. It is a layered, multi-service, distributed lakehouse platform with explicit operational logic and clear data movement boundaries. On this basis, the next chapter can focus on deployment strategy and distributed execution perspective, explaining how this implementation is physically organized and operated across the final multi-node environment.
